\section{账本事件的社会制度深描：early-ming}

\label{sec:ledger-depth-early-ming}

以下各节依本卷账本的事件、来源和置信提示组织，不以编年节点替代地方社会。朝廷命令、法律条文、城市记录、家庭文书与物质遗存的证明范围不同，必须让其差异保留在叙事中。

\subsection{政治制度与地方执行}

在账本条目\texttt{undefined}中，1368（洪武元年）八月的“徐达军进入大都并改称北平”发生于明与北元。本条把背景说明为明军自山东、河南北上且元顺帝已北撤；大都守备难获外援，并以元廷失去旧都和华北行政中心；明须经营北平并面对草原残余政权概括可见影响。要理解它的社会后果，不能止于宫廷或战场：土地、赋役、司法、道路、性别化家庭劳动、城市市场、书院寺观和边海网络都会影响地方如何接收或改写制度。

本条依据《太祖实录》洪武元年八月条；《明史·太祖本纪》；《高丽史·恭愍王世家》二十一年条。这些来源能够较稳妥地支持所记年序、诏令、任命或特定地点的行动，但无法自动证明全国的财产、人口、识字、信仰或动机。账本保留的证据说明是“日期可靠；明廷军报不独证城内损失和居民去留”。所以，官府标签、奏疏数字与后来的道德评语必须放回其文书目的，不能被写成不经分区核验的社会全貌。

这条事件更适合作为一组关系变化的锚点：中央方案需由地方官、书吏、士绅、宗族、商人与劳动者共同转译；不同家庭面对的税役、婚姻、迁徙和安全风险并不相等。材料只让我们看见其中有限的记录者。承认可见与不可见的界线，才能使制度史、文化史和区域史互相校正。

在账本条目\texttt{undefined}中，1368（洪武元年）正月的“朱元璋在应天即皇帝位并建国号明”发生于明。本条把背景说明为红巾军政权分裂而元廷南方控制崩解；朱元璋已控制江南粮赋与应天军政中心，并以明廷取得颁诏和任官名义；北伐与地方接收须转为连续的行政和军事行动概括可见影响。要理解它的社会后果，不能止于宫廷或战场：土地、赋役、司法、道路、性别化家庭劳动、城市市场、书院寺观和边海网络都会影响地方如何接收或改写制度。

本条依据《太祖实录》洪武元年正月条（即位诏）；《明史·太祖本纪》；南京明故宫遗址考古发掘报告。这些来源能够较稳妥地支持所记年序、诏令、任命或特定地点的行动，但无法自动证明全国的财产、人口、识字、信仰或动机。账本保留的证据说明是“日期可靠；诏告可证建国程序而不能替代各地归附过程”。所以，官府标签、奏疏数字与后来的道德评语必须放回其文书目的，不能被写成不经分区核验的社会全貌。

这条事件更适合作为一组关系变化的锚点：中央方案需由地方官、书吏、士绅、宗族、商人与劳动者共同转译；不同家庭面对的税役、婚姻、迁徙和安全风险并不相等。材料只让我们看见其中有限的记录者。承认可见与不可见的界线，才能使制度史、文化史和区域史互相校正。

在账本条目\texttt{undefined}中，1368（洪武元年）九月的“在南京设国子学以培训中央官员子弟”发生于明。本条把背景说明为新朝须为六部和地方官体系补充识字行政人才并重建礼制教育，并以南京国子监成为明初官学枢纽；科举与荐举的关系仍待后续制度调整概括可见影响。要理解它的社会后果，不能止于宫廷或战场：土地、赋役、司法、道路、性别化家庭劳动、城市市场、书院寺观和边海网络都会影响地方如何接收或改写制度。

本条依据《太祖实录》洪武元年九月条；《明会典·国子监》；《明史·选举志》。这些来源能够较稳妥地支持所记年序、诏令、任命或特定地点的行动，但无法自动证明全国的财产、人口、识字、信仰或动机。账本保留的证据说明是“日期较可靠；学校建置不等于生员规模或教学成效”。所以，官府标签、奏疏数字与后来的道德评语必须放回其文书目的，不能被写成不经分区核验的社会全貌。

这条事件更适合作为一组关系变化的锚点：中央方案需由地方官、书吏、士绅、宗族、商人与劳动者共同转译；不同家庭面对的税役、婚姻、迁徙和安全风险并不相等。材料只让我们看见其中有限的记录者。承认可见与不可见的界线，才能使制度史、文化史和区域史互相校正。

在账本条目\texttt{undefined}中，1369（洪武二年）五月至八月的“明军北进上都、应昌一线；元顺帝在北遁途中去世”发生于明与北元。本条把背景说明为大都失守后北元仍试图依托上都、应昌和草原骑兵维持王朝中心，并以明军未能消灭草原政治力量；北边战争由夺城转为长期机动与互市问题概括可见影响。要理解它的社会后果，不能止于宫廷或战场：土地、赋役、司法、道路、性别化家庭劳动、城市市场、书院寺观和边海网络都会影响地方如何接收或改写制度。

本条依据《太祖实录》洪武二年五月、八月条；《明史·北虏传》；《高丽史》北元消息条。这些来源能够较稳妥地支持所记年序、诏令、任命或特定地点的行动，但无法自动证明全国的财产、人口、识字、信仰或动机。账本保留的证据说明是“日期与战事序列较可靠；元顺帝死地和撤退细节有异说”。所以，官府标签、奏疏数字与后来的道德评语必须放回其文书目的，不能被写成不经分区核验的社会全貌。

这条事件更适合作为一组关系变化的锚点：中央方案需由地方官、书吏、士绅、宗族、商人与劳动者共同转译；不同家庭面对的税役、婚姻、迁徙和安全风险并不相等。材料只让我们看见其中有限的记录者。承认可见与不可见的界线，才能使制度史、文化史和区域史互相校正。

在账本条目\texttt{undefined}中，1370（洪武三年）四月的“太祖大封诸子为王并配置藩国”发生于明。本条把背景说明为建国后皇室需要可继承的资源分配与边地防卫网络，并以诸王成为地方军政力量；中央与藩王的权限张力为建文时期冲突埋下制度条件概括可见影响。要理解它的社会后果，不能止于宫廷或战场：土地、赋役、司法、道路、性别化家庭劳动、城市市场、书院寺观和边海网络都会影响地方如何接收或改写制度。

本条依据《太祖实录》洪武三年四月条；《明史·诸王传》；《明会典·宗藩》。这些来源能够较稳妥地支持所记年序、诏令、任命或特定地点的行动，但无法自动证明全国的财产、人口、识字、信仰或动机。账本保留的证据说明是“日期可靠；封册可证名义授封而不能直接证明各王实际兵权”。所以，官府标签、奏疏数字与后来的道德评语必须放回其文书目的，不能被写成不经分区核验的社会全貌。

这条事件更适合作为一组关系变化的锚点：中央方案需由地方官、书吏、士绅、宗族、商人与劳动者共同转译；不同家庭面对的税役、婚姻、迁徙和安全风险并不相等。材料只让我们看见其中有限的记录者。承认可见与不可见的界线，才能使制度史、文化史和区域史互相校正。

在账本条目\texttt{undefined}中，1370（洪武三年）六月的“徐达在沈儿峪击败扩廓帖木儿部”发生于明与北元。本条把背景说明为明廷欲趁北元重组未定追击其主力；徐达军依托山西、陕西粮道北进，并以扩廓部受挫但未被歼灭；明廷随后仍需以多路远征争夺草原主动权概括可见影响。要理解它的社会后果，不能止于宫廷或战场：土地、赋役、司法、道路、性别化家庭劳动、城市市场、书院寺观和边海网络都会影响地方如何接收或改写制度。

本条依据《太祖实录》洪武三年六月条；《明史·徐达传》与《北虏传》；17世纪《蒙古源流》北元汗统叙事。这些来源能够较稳妥地支持所记年序、诏令、任命或特定地点的行动，但无法自动证明全国的财产、人口、识字、信仰或动机。账本保留的证据说明是“战役发生可靠；双方兵数和斩获主要来自明方战报”。所以，官府标签、奏疏数字与后来的道德评语必须放回其文书目的，不能被写成不经分区核验的社会全貌。

这条事件更适合作为一组关系变化的锚点：中央方案需由地方官、书吏、士绅、宗族、商人与劳动者共同转译；不同家庭面对的税役、婚姻、迁徙和安全风险并不相等。材料只让我们看见其中有限的记录者。承认可见与不可见的界线，才能使制度史、文化史和区域史互相校正。

在账本条目\texttt{undefined}中，1371（洪武四年）十二月的“朝廷申禁滨海居民私自出海并限制民间海外贸易”发生于明与东南沿海社会。本条把背景说明为倭患警报、元末海上武装遗绪与新朝对沿海人口流动的疑虑交织，并以朝贡贸易被置于官府控制之下；私航、走私和沿海防务成为持续的地方—中央博弈概括可见影响。要理解它的社会后果，不能止于宫廷或战场：土地、赋役、司法、道路、性别化家庭劳动、城市市场、书院寺观和边海网络都会影响地方如何接收或改写制度。

本条依据《太祖实录》洪武四年十二月条；《明史·食货志》；《筹海图编》所录早明海禁材料。这些来源能够较稳妥地支持所记年序、诏令、任命或特定地点的行动，但无法自动证明全国的财产、人口、识字、信仰或动机。账本保留的证据说明是“诏令日期可靠；禁令不等于沿海航行和交易立即停止”。所以，官府标签、奏疏数字与后来的道德评语必须放回其文书目的，不能被写成不经分区核验的社会全貌。

这条事件更适合作为一组关系变化的锚点：中央方案需由地方官、书吏、士绅、宗族、商人与劳动者共同转译；不同家庭面对的税役、婚姻、迁徙和安全风险并不相等。材料只让我们看见其中有限的记录者。承认可见与不可见的界线，才能使制度史、文化史和区域史互相校正。

在账本条目\texttt{undefined}中，1371（洪武四年）的“傅友德等沿长江入蜀，夏主明昇降明”发生于明与明夏。本条把背景说明为明夏控制四川但与中原补给和政治盟友隔绝；明军掌握长江中下游军运条件，并以四川纳入明朝行政与赋役框架；山地边区和土司关系仍须逐步处理概括可见影响。要理解它的社会后果，不能止于宫廷或战场：土地、赋役、司法、道路、性别化家庭劳动、城市市场、书院寺观和边海网络都会影响地方如何接收或改写制度。

本条依据《太祖实录》洪武四年条；《明史·明玉珍传》；《四川通志》明夏遗迹条。这些来源能够较稳妥地支持所记年序、诏令、任命或特定地点的行动，但无法自动证明全国的财产、人口、识字、信仰或动机。账本保留的证据说明是“年代可靠；行军路线和地方响应不能只据凯旋叙事复原”。所以，官府标签、奏疏数字与后来的道德评语必须放回其文书目的，不能被写成不经分区核验的社会全貌。

这条事件更适合作为一组关系变化的锚点：中央方案需由地方官、书吏、士绅、宗族、商人与劳动者共同转译；不同家庭面对的税役、婚姻、迁徙和安全风险并不相等。材料只让我们看见其中有限的记录者。承认可见与不可见的界线，才能使制度史、文化史和区域史互相校正。

在账本条目\texttt{undefined}中，1372（洪武五年）的“三路北伐受挫，李文忠、徐达等部未能取得预期战果”发生于明与北元。本条把背景说明为明廷试图以多路大军压迫北元；行军距离、草原水草和部队协同限制远征，并以太祖转向更审慎的北边经营；卫所、边墙和藩王防线的重要性上升概括可见影响。要理解它的社会后果，不能止于宫廷或战场：土地、赋役、司法、道路、性别化家庭劳动、城市市场、书院寺观和边海网络都会影响地方如何接收或改写制度。

本条依据《太祖实录》洪武五年条；《明史·李文忠传》与《北虏传》；17世纪《蒙古源流》北元汗统叙事。这些来源能够较稳妥地支持所记年序、诏令、任命或特定地点的行动，但无法自动证明全国的财产、人口、识字、信仰或动机。账本保留的证据说明是“战役大势可靠；败因、损失和各路会师情况在记载中不一”。所以，官府标签、奏疏数字与后来的道德评语必须放回其文书目的，不能被写成不经分区核验的社会全貌。

这条事件更适合作为一组关系变化的锚点：中央方案需由地方官、书吏、士绅、宗族、商人与劳动者共同转译；不同家庭面对的税役、婚姻、迁徙和安全风险并不相等。材料只让我们看见其中有限的记录者。承认可见与不可见的界线，才能使制度史、文化史和区域史互相校正。

在账本条目\texttt{undefined}中，1372（洪武五年）的“明使杨载出使琉球并开启朝贡往来”发生于明与琉球诸政权。本条把背景说明为海禁后明廷仍需以官方使节组织可控的海上外交；琉球各势力寻求外部承认，并以形成持续的册封—贸易通道；朝贡名义与实际航运、商人网络不能混同概括可见影响。要理解它的社会后果，不能止于宫廷或战场：土地、赋役、司法、道路、性别化家庭劳动、城市市场、书院寺观和边海网络都会影响地方如何接收或改写制度。

本条依据《太祖实录》洪武五年条；《明史·琉球传》；《历代宝案》早期入贡文书。这些来源能够较稳妥地支持所记年序、诏令、任命或特定地点的行动，但无法自动证明全国的财产、人口、识字、信仰或动机。账本保留的证据说明是“明使出航与琉球回应有双边记录；三山内部政治不应视作单一国家”。所以，官府标签、奏疏数字与后来的道德评语必须放回其文书目的，不能被写成不经分区核验的社会全貌。

这条事件更适合作为一组关系变化的锚点：中央方案需由地方官、书吏、士绅、宗族、商人与劳动者共同转译；不同家庭面对的税役、婚姻、迁徙和安全风险并不相等。材料只让我们看见其中有限的记录者。承认可见与不可见的界线，才能使制度史、文化史和区域史互相校正。

在账本条目\texttt{undefined}中，1373（洪武六年）的“太祖停罢科举，改重荐举与学校考核”发生于明。本条把背景说明为首次科举所取士人与新朝行政需求之间发生张力；皇帝偏好直接控制人才来源，并以荐举和学校渠道暂时扩张；科举何时、如何恢复成为中央选官制度的后续问题概括可见影响。要理解它的社会后果，不能止于宫廷或战场：土地、赋役、司法、道路、性别化家庭劳动、城市市场、书院寺观和边海网络都会影响地方如何接收或改写制度。

本条依据《太祖实录》洪武六年条；《明史·选举志》；《明会典·礼部》科举沿革。这些来源能够较稳妥地支持所记年序、诏令、任命或特定地点的行动，但无法自动证明全国的财产、人口、识字、信仰或动机。账本保留的证据说明是“政策日期较可靠；官方对取士质量的评语带有皇帝立场”。所以，官府标签、奏疏数字与后来的道德评语必须放回其文书目的，不能被写成不经分区核验的社会全貌。

这条事件更适合作为一组关系变化的锚点：中央方案需由地方官、书吏、士绅、宗族、商人与劳动者共同转译；不同家庭面对的税役、婚姻、迁徙和安全风险并不相等。材料只让我们看见其中有限的记录者。承认可见与不可见的界线，才能使制度史、文化史和区域史互相校正。

\subsection{法律、家庭与社会关系}

在账本条目\texttt{undefined}中，1376（洪武九年）六月的“废行中书省，分置承宣布政使司、都指挥使司、提刑按察使司”发生于明。本条把背景说明为太祖防范行省集军政财权于一体，并需将新征服地区纳入可分察的行政链条，并以地方行政、军事和监察被分别归属三司；协调成本和中央派遣机制随之增加概括可见影响。要理解它的社会后果，不能止于宫廷或战场：土地、赋役、司法、道路、性别化家庭劳动、城市市场、书院寺观和边海网络都会影响地方如何接收或改写制度。

本条依据《太祖实录》洪武九年六月条；《明史·职官志》；《明会典·布政司》。这些来源能够较稳妥地支持所记年序、诏令、任命或特定地点的行动，但无法自动证明全国的财产、人口、识字、信仰或动机。账本保留的证据说明是“制度发布可靠；三司分置不等于地方权力已均衡运作”。所以，官府标签、奏疏数字与后来的道德评语必须放回其文书目的，不能被写成不经分区核验的社会全貌。

这条事件更适合作为一组关系变化的锚点：中央方案需由地方官、书吏、士绅、宗族、商人与劳动者共同转译；不同家庭面对的税役、婚姻、迁徙和安全风险并不相等。材料只让我们看见其中有限的记录者。承认可见与不可见的界线，才能使制度史、文化史和区域史互相校正。

在账本条目\texttt{undefined}中，1380（洪武十三年）正月的“胡惟庸案后中书省被废，皇帝直接统辖六部”发生于明。本条把背景说明为太祖对宰相权力和官僚结党高度警惕；案件成为重组中枢权力的触发点，并以六部直接对皇帝负责；政务依赖皇帝裁断与后续内阁文书协调的矛盾加深概括可见影响。要理解它的社会后果，不能止于宫廷或战场：土地、赋役、司法、道路、性别化家庭劳动、城市市场、书院寺观和边海网络都会影响地方如何接收或改写制度。

本条依据《太祖实录》洪武十三年正月条；《明史·胡惟庸传》与《职官志》；《大诰》胡惟庸案条。这些来源能够较稳妥地支持所记年序、诏令、任命或特定地点的行动，但无法自动证明全国的财产、人口、识字、信仰或动机。账本保留的证据说明是“废相与制度变动可靠；胡惟庸“通倭通元”等指控须视为案狱叙事而非既定事实”。所以，官府标签、奏疏数字与后来的道德评语必须放回其文书目的，不能被写成不经分区核验的社会全貌。

这条事件更适合作为一组关系变化的锚点：中央方案需由地方官、书吏、士绅、宗族、商人与劳动者共同转译；不同家庭面对的税役、婚姻、迁徙和安全风险并不相等。材料只让我们看见其中有限的记录者。承认可见与不可见的界线，才能使制度史、文化史和区域史互相校正。

在账本条目\texttt{undefined}中，1381（洪武十四年）的“朝廷推进黄册与里甲编制，登记户口、田土和役属”发生于明与里甲户籍社会。本条把背景说明为新朝须将赋役、军户和土地责任落实到可追查的户等；元末流亡和战争破坏使登记困难，并以里甲成为赋役征派的基层框架；漏报、逃役和地方中介将持续塑造实际负担概括可见影响。要理解它的社会后果，不能止于宫廷或战场：土地、赋役、司法、道路、性别化家庭劳动、城市市场、书院寺观和边海网络都会影响地方如何接收或改写制度。

本条依据《太祖实录》洪武十四年户籍条；《明会典·户口》；徽州等地明代黄册残卷。这些来源能够较稳妥地支持所记年序、诏令、任命或特定地点的行动，但无法自动证明全国的财产、人口、识字、信仰或动机。账本保留的证据说明是“制度命令可靠；存世黄册仅反映部分地区与登记口径”。所以，官府标签、奏疏数字与后来的道德评语必须放回其文书目的，不能被写成不经分区核验的社会全貌。

这条事件更适合作为一组关系变化的锚点：中央方案需由地方官、书吏、士绅、宗族、商人与劳动者共同转译；不同家庭面对的税役、婚姻、迁徙和安全风险并不相等。材料只让我们看见其中有限的记录者。承认可见与不可见的界线，才能使制度史、文化史和区域史互相校正。

在账本条目\texttt{undefined}中，1381（洪武十四年）的“傅友德、蓝玉、沐英率军自川黔进攻云南”发生于明、梁王与云南诸势力。本条把背景说明为梁王仍据云南并与北元保持联系；明廷欲控制西南交通、矿产和边疆名义，并以大规模军事接收开始；军屯、卫所和土司安排将深刻改变云南地方秩序概括可见影响。要理解它的社会后果，不能止于宫廷或战场：土地、赋役、司法、道路、性别化家庭劳动、城市市场、书院寺观和边海网络都会影响地方如何接收或改写制度。

本条依据《太祖实录》洪武十四年条；《明史·傅友德传》《沐英传》；《云南图经志书》云南府条。这些来源能够较稳妥地支持所记年序、诏令、任命或特定地点的行动，但无法自动证明全国的财产、人口、识字、信仰或动机。账本保留的证据说明是“出兵和主帅可靠；军额、伤亡与各族群立场不可照录明方奏报”。所以，官府标签、奏疏数字与后来的道德评语必须放回其文书目的，不能被写成不经分区核验的社会全貌。

这条事件更适合作为一组关系变化的锚点：中央方案需由地方官、书吏、士绅、宗族、商人与劳动者共同转译；不同家庭面对的税役、婚姻、迁徙和安全风险并不相等。材料只让我们看见其中有限的记录者。承认可见与不可见的界线，才能使制度史、文化史和区域史互相校正。

在账本条目\texttt{undefined}中，1382（洪武十五年）的“昆明陷落后梁王自尽，明廷在云南设卫所并逐步建置”发生于明与云南地方政权。本条把背景说明为明军取得昆明后需维持补给并处理旧元官军、土司和山地社群，并以云南成为军屯和土司并行的边疆区域；征调成本与地方反抗未因占城而结束概括可见影响。要理解它的社会后果，不能止于宫廷或战场：土地、赋役、司法、道路、性别化家庭劳动、城市市场、书院寺观和边海网络都会影响地方如何接收或改写制度。

本条依据《太祖实录》洪武十五年条；《明史·云南土司传》；大理、昆明地区碑刻与地方志。这些来源能够较稳妥地支持所记年序、诏令、任命或特定地点的行动，但无法自动证明全国的财产、人口、识字、信仰或动机。账本保留的证据说明是“核心结果可靠；地方降附过程和破坏程度需以地方材料分别核验”。所以，官府标签、奏疏数字与后来的道德评语必须放回其文书目的，不能被写成不经分区核验的社会全貌。

这条事件更适合作为一组关系变化的锚点：中央方案需由地方官、书吏、士绅、宗族、商人与劳动者共同转译；不同家庭面对的税役、婚姻、迁徙和安全风险并不相等。材料只让我们看见其中有限的记录者。承认可见与不可见的界线，才能使制度史、文化史和区域史互相校正。

在账本条目\texttt{undefined}中，1383（洪武十六年）的“设置锦衣卫，皇帝以亲军兼掌侦缉和诏狱”发生于明。本条把背景说明为废相后皇帝强化对官僚和军政信息的直接掌握；大案审讯需要绕开常规司法链，并以亲军侦缉成为中枢权力工具之一；司法程序与政治案件的界线更趋紧张概括可见影响。要理解它的社会后果，不能止于宫廷或战场：土地、赋役、司法、道路、性别化家庭劳动、城市市场、书院寺观和边海网络都会影响地方如何接收或改写制度。

本条依据《太祖实录》洪武十六年条；《明史·职官志》；《明会典·锦衣卫》。这些来源能够较稳妥地支持所记年序、诏令、任命或特定地点的行动，但无法自动证明全国的财产、人口、识字、信仰或动机。账本保留的证据说明是“设卫可靠；后世将其职能倒投到洪武初年须谨慎”。所以，官府标签、奏疏数字与后来的道德评语必须放回其文书目的，不能被写成不经分区核验的社会全貌。

这条事件更适合作为一组关系变化的锚点：中央方案需由地方官、书吏、士绅、宗族、商人与劳动者共同转译；不同家庭面对的税役、婚姻、迁徙和安全风险并不相等。材料只让我们看见其中有限的记录者。承认可见与不可见的界线，才能使制度史、文化史和区域史互相校正。

在账本条目\texttt{undefined}中，1385（洪武十八年）的“《大诰》颁行，以案例和重刑语言约束官民”发生于明。本条把背景说明为太祖认为常律不足以威慑吏治和社会流动中的违法行为，并借案例塑造服从语言，并以法外重典进入基层教化与审判语境；法条、诏令和地方执行之间的落差须另行追踪概括可见影响。要理解它的社会后果，不能止于宫廷或战场：土地、赋役、司法、道路、性别化家庭劳动、城市市场、书院寺观和边海网络都会影响地方如何接收或改写制度。

本条依据《大诰》初编；《太祖实录》洪武十八年条；《明史·刑法志》。这些来源能够较稳妥地支持所记年序、诏令、任命或特定地点的行动，但无法自动证明全国的财产、人口、识字、信仰或动机。账本保留的证据说明是“文本与颁行可靠；规范性惩戒话语不能直接换算为实际犯罪率”。所以，官府标签、奏疏数字与后来的道德评语必须放回其文书目的，不能被写成不经分区核验的社会全貌。

这条事件更适合作为一组关系变化的锚点：中央方案需由地方官、书吏、士绅、宗族、商人与劳动者共同转译；不同家庭面对的税役、婚姻、迁徙和安全风险并不相等。材料只让我们看见其中有限的记录者。承认可见与不可见的界线，才能使制度史、文化史和区域史互相校正。

在账本条目\texttt{undefined}中，1385（洪武十八年）的“郭桓案牵连户部及各地征粮官吏，朝廷严惩侵没赋税者”发生于明与户部、地方官吏。本条把背景说明为江南赋粮征解经多层胥吏和官员转运，中央难以核对定额、耗损与实收，并以恐惧性稽核整顿财政链条但也冲击地方行政；账册真实度与官吏自保问题加剧概括可见影响。要理解它的社会后果，不能止于宫廷或战场：土地、赋役、司法、道路、性别化家庭劳动、城市市场、书院寺观和边海网络都会影响地方如何接收或改写制度。

本条依据《太祖实录》洪武十八年条；《明史·刑法志》；《大诰》郭桓案条。这些来源能够较稳妥地支持所记年序、诏令、任命或特定地点的行动，但无法自动证明全国的财产、人口、识字、信仰或动机。账本保留的证据说明是“案件与惩处可靠；被杀人数及“全国贪墨”规模高度依赖官方案狱口径”。所以，官府标签、奏疏数字与后来的道德评语必须放回其文书目的，不能被写成不经分区核验的社会全貌。

这条事件更适合作为一组关系变化的锚点：中央方案需由地方官、书吏、士绅、宗族、商人与劳动者共同转译；不同家庭面对的税役、婚姻、迁徙和安全风险并不相等。材料只让我们看见其中有限的记录者。承认可见与不可见的界线，才能使制度史、文化史和区域史互相校正。

在账本条目\texttt{undefined}中，1387（洪武二十年）的“朝廷催成鱼鳞图册等田土登记以配合赋役核算”发生于明与江南田土社会。本条把背景说明为黄册登记户等后，国家需以地块和赋额关联田土责任；土地转移与隐匿造成核算压力，并以田土登记成为赋役分配的凭据；地方书手、业主和契约实践影响其可执行性概括可见影响。要理解它的社会后果，不能止于宫廷或战场：土地、赋役、司法、道路、性别化家庭劳动、城市市场、书院寺观和边海网络都会影响地方如何接收或改写制度。

本条依据《太祖实录》洪武二十年田土条；《明会典·田土》；徽州鱼鳞图册与相关契约。这些来源能够较稳妥地支持所记年序、诏令、任命或特定地点的行动，但无法自动证明全国的财产、人口、识字、信仰或动机。账本保留的证据说明是“命令与部分册籍可证；图册年代、抄补层次和亩数不等于实际占有”。所以，官府标签、奏疏数字与后来的道德评语必须放回其文书目的，不能被写成不经分区核验的社会全貌。

这条事件更适合作为一组关系变化的锚点：中央方案需由地方官、书吏、士绅、宗族、商人与劳动者共同转译；不同家庭面对的税役、婚姻、迁徙和安全风险并不相等。材料只让我们看见其中有限的记录者。承认可见与不可见的界线，才能使制度史、文化史和区域史互相校正。

在账本条目\texttt{undefined}中，1387（洪武二十年）的“冯胜等在辽东、松花江方向迫使纳哈出率部降明”发生于明与纳哈出部。本条把背景说明为北元残部在辽东与松花江流域保持机动空间；明廷需稳住山海关以东通道，并以明在辽东增强卫所布置与招抚能力；东北诸部的朝贡和军事关系仍不断变动概括可见影响。要理解它的社会后果，不能止于宫廷或战场：土地、赋役、司法、道路、性别化家庭劳动、城市市场、书院寺观和边海网络都会影响地方如何接收或改写制度。

本条依据《太祖实录》洪武二十年条；《明史·冯胜传》与《北虏传》；《高丽史》辽东方面条。这些来源能够较稳妥地支持所记年序、诏令、任命或特定地点的行动，但无法自动证明全国的财产、人口、识字、信仰或动机。账本保留的证据说明是“降附事实可靠；部众数量和“归化”程度不能由明方献俘叙事推定”。所以，官府标签、奏疏数字与后来的道德评语必须放回其文书目的，不能被写成不经分区核验的社会全貌。

这条事件更适合作为一组关系变化的锚点：中央方案需由地方官、书吏、士绅、宗族、商人与劳动者共同转译；不同家庭面对的税役、婚姻、迁徙和安全风险并不相等。材料只让我们看见其中有限的记录者。承认可见与不可见的界线，才能使制度史、文化史和区域史互相校正。

在账本条目\texttt{undefined}中，1388（洪武二十一年）的“蓝玉军在捕鱼儿海击溃北元主力并俘获大批人员物资”发生于明与北元。本条把背景说明为明军利用北元内部不稳和长距离奔袭寻求决定性战果，并以北元汗庭进一步分散；蓝玉军功骤升，也加深洪武朝对勋贵兵权的疑惧概括可见影响。要理解它的社会后果，不能止于宫廷或战场：土地、赋役、司法、道路、性别化家庭劳动、城市市场、书院寺观和边海网络都会影响地方如何接收或改写制度。

本条依据《太祖实录》洪武二十一年条；《明史·蓝玉传》与《北虏传》；17世纪《蒙古源流》北元汗统叙事。这些来源能够较稳妥地支持所记年序、诏令、任命或特定地点的行动，但无法自动证明全国的财产、人口、识字、信仰或动机。账本保留的证据说明是“胜利可靠；战果数字、元主脱逃经过和俘获规模应保留两方材料差异”。所以，官府标签、奏疏数字与后来的道德评语必须放回其文书目的，不能被写成不经分区核验的社会全貌。

这条事件更适合作为一组关系变化的锚点：中央方案需由地方官、书吏、士绅、宗族、商人与劳动者共同转译；不同家庭面对的税役、婚姻、迁徙和安全风险并不相等。材料只让我们看见其中有限的记录者。承认可见与不可见的界线，才能使制度史、文化史和区域史互相校正。

\subsection{城市、文化与知识流通}

在账本条目\texttt{undefined}中，1390（洪武二十三年）的“蓝玉以谋反罪被诛，功臣集团再受清洗”发生于明。本条把背景说明为北伐功臣掌握军中网络且太祖持续防范结党；案狱成为重排军政人事的渠道，并以开国武将集团受到重创；皇太孙继承安排更依赖文官和制度性控制概括可见影响。要理解它的社会后果，不能止于宫廷或战场：土地、赋役、司法、道路、性别化家庭劳动、城市市场、书院寺观和边海网络都会影响地方如何接收或改写制度。

本条依据《太祖实录》洪武二十三年条；《明史·蓝玉传》；《大诰》相关案件。这些来源能够较稳妥地支持所记年序、诏令、任命或特定地点的行动，但无法自动证明全国的财产、人口、识字、信仰或动机。账本保留的证据说明是“处决与案名可靠；谋反细节主要来自后修实录和官方传记，难以独立裁定”。所以，官府标签、奏疏数字与后来的道德评语必须放回其文书目的，不能被写成不经分区核验的社会全貌。

这条事件更适合作为一组关系变化的锚点：中央方案需由地方官、书吏、士绅、宗族、商人与劳动者共同转译；不同家庭面对的税役、婚姻、迁徙和安全风险并不相等。材料只让我们看见其中有限的记录者。承认可见与不可见的界线，才能使制度史、文化史和区域史互相校正。

在账本条目\texttt{undefined}中，1392（洪武二十五年）四月的“皇太子朱标去世，太祖立朱允炆为皇太孙”发生于明。本条把背景说明为朱标早逝打破原有继承顺序；太祖需在皇孙与成年藩王之间确立可被承认的名分，并以建文君臣日后削藩拥有继承政治背景；燕王等边王的资源与疑虑同时上升概括可见影响。要理解它的社会后果，不能止于宫廷或战场：土地、赋役、司法、道路、性别化家庭劳动、城市市场、书院寺观和边海网络都会影响地方如何接收或改写制度。

本条依据《太祖实录》洪武二十五年四月条；《明史·懿文太子传》与《惠帝本纪》；南京懿文太子陵考古简报与碑刻著录。这些来源能够较稳妥地支持所记年序、诏令、任命或特定地点的行动，但无法自动证明全国的财产、人口、识字、信仰或动机。账本保留的证据说明是“继承变更可靠；对朱标能力和诸王态度的评价多为后出叙事”。所以，官府标签、奏疏数字与后来的道德评语必须放回其文书目的，不能被写成不经分区核验的社会全貌。

这条事件更适合作为一组关系变化的锚点：中央方案需由地方官、书吏、士绅、宗族、商人与劳动者共同转译；不同家庭面对的税役、婚姻、迁徙和安全风险并不相等。材料只让我们看见其中有限的记录者。承认可见与不可见的界线，才能使制度史、文化史和区域史互相校正。

在账本条目\texttt{undefined}中，1393（洪武二十六年）的“蓝玉案牵连扩大，勋贵与官员大批获罪”发生于明。本条把背景说明为太祖以案狱消除潜在军事和官僚自主网络，并为皇太孙预作权力清场，并以朝廷人事更加依附皇权；以高压保障继承也削弱了可制衡藩王的开国武臣概括可见影响。要理解它的社会后果，不能止于宫廷或战场：土地、赋役、司法、道路、性别化家庭劳动、城市市场、书院寺观和边海网络都会影响地方如何接收或改写制度。

本条依据《太祖实录》洪武二十六年条；《明史·功臣世表》与相关列传；《大诰三编》。这些来源能够较稳妥地支持所记年序、诏令、任命或特定地点的行动，但无法自动证明全国的财产、人口、识字、信仰或动机。账本保留的证据说明是“清洗持续及其政治影响可靠；受牵连人数和罪证须避免采纳单一官方总数”。所以，官府标签、奏疏数字与后来的道德评语必须放回其文书目的，不能被写成不经分区核验的社会全貌。

这条事件更适合作为一组关系变化的锚点：中央方案需由地方官、书吏、士绅、宗族、商人与劳动者共同转译；不同家庭面对的税役、婚姻、迁徙和安全风险并不相等。材料只让我们看见其中有限的记录者。承认可见与不可见的界线，才能使制度史、文化史和区域史互相校正。

在账本条目\texttt{undefined}中，1397（洪武三十年）的“定颁《大明律》终本，形成以六律为纲的成文法框架”发生于明。本条把背景说明为洪武前期反复修律与《大诰》重典并行，中央需要稳定、可传抄的法典文本，并以终本成为明代司法基本参照；诏令、条例和诏狱仍会改变实际司法面貌概括可见影响。要理解它的社会后果，不能止于宫廷或战场：土地、赋役、司法、道路、性别化家庭劳动、城市市场、书院寺观和边海网络都会影响地方如何接收或改写制度。

本条依据《大明律》洪武三十年刊本系统；《太祖实录》洪武三十年条；《明史·刑法志》。这些来源能够较稳妥地支持所记年序、诏令、任命或特定地点的行动，但无法自动证明全国的财产、人口、识字、信仰或动机。账本保留的证据说明是“终本年代可靠；法典条文不能证明州县审判或社会遵从一致”。所以，官府标签、奏疏数字与后来的道德评语必须放回其文书目的，不能被写成不经分区核验的社会全貌。

这条事件更适合作为一组关系变化的锚点：中央方案需由地方官、书吏、士绅、宗族、商人与劳动者共同转译；不同家庭面对的税役、婚姻、迁徙和安全风险并不相等。材料只让我们看见其中有限的记录者。承认可见与不可见的界线，才能使制度史、文化史和区域史互相校正。

在账本条目\texttt{undefined}中，1397（洪武三十年）的“会试南北榜案后重取中式者，籍贯分布被纳入取士政治”发生于明。本条把背景说明为江南考生在考试中占多，北方士人及皇帝对选才地域失衡产生质疑，并以取士过程被强化为皇权可干预的政治问题；南北地域名额与资格争论延续概括可见影响。要理解它的社会后果，不能止于宫廷或战场：土地、赋役、司法、道路、性别化家庭劳动、城市市场、书院寺观和边海网络都会影响地方如何接收或改写制度。

本条依据《太祖实录》洪武三十年条；《明史·选举志》；《明会典·礼部》科举条。这些来源能够较稳妥地支持所记年序、诏令、任命或特定地点的行动，但无法自动证明全国的财产、人口、识字、信仰或动机。账本保留的证据说明是“事件可靠；“南北地域偏私”是朝廷处理语言，不能还原为单一真实原因”。所以，官府标签、奏疏数字与后来的道德评语必须放回其文书目的，不能被写成不经分区核验的社会全貌。

这条事件更适合作为一组关系变化的锚点：中央方案需由地方官、书吏、士绅、宗族、商人与劳动者共同转译；不同家庭面对的税役、婚姻、迁徙和安全风险并不相等。材料只让我们看见其中有限的记录者。承认可见与不可见的界线，才能使制度史、文化史和区域史互相校正。

在账本条目\texttt{undefined}中，1398（洪武三十一年）闰五月的“太祖去世，皇太孙朱允炆即位为建文帝”发生于明。本条把背景说明为太祖已以皇太孙继承并清洗部分功臣，但成年藩王仍掌边镇卫所资源，并以建文新政面对削藩与守成的选择；中央—燕王冲突迅速制度化概括可见影响。要理解它的社会后果，不能止于宫廷或战场：土地、赋役、司法、道路、性别化家庭劳动、城市市场、书院寺观和边海网络都会影响地方如何接收或改写制度。

本条依据《太祖实录》洪武三十一年闰五月条；《明史·惠帝本纪》；孝陵碑刻与陵寝考古资料。这些来源能够较稳妥地支持所记年序、诏令、任命或特定地点的行动，但无法自动证明全国的财产、人口、识字、信仰或动机。账本保留的证据说明是“继位日期可靠；遗诏与宫廷部署的细节受建文、永乐两朝修史影响”。所以，官府标签、奏疏数字与后来的道德评语必须放回其文书目的，不能被写成不经分区核验的社会全貌。

这条事件更适合作为一组关系变化的锚点：中央方案需由地方官、书吏、士绅、宗族、商人与劳动者共同转译；不同家庭面对的税役、婚姻、迁徙和安全风险并不相等。材料只让我们看见其中有限的记录者。承认可见与不可见的界线，才能使制度史、文化史和区域史互相校正。

在账本条目\texttt{undefined}中，1399（建文元年）七月的“燕王朱棣以“靖难”为名起兵反对建文朝”发生于明与燕王集团。本条把背景说明为建文朝先后处理周、齐、湘等王，燕王担忧自身兵权和封国安全，并以内战使北方卫所、运河和山东交通成为争夺对象；胜负将决定继承叙事的书写者概括可见影响。要理解它的社会后果，不能止于宫廷或战场：土地、赋役、司法、道路、性别化家庭劳动、城市市场、书院寺观和边海网络都会影响地方如何接收或改写制度。

本条依据《明太宗实录》建文元年七月条；《明史·成祖本纪》与《方孝孺传》；北平地方碑刻和方志线索。这些来源能够较稳妥地支持所记年序、诏令、任命或特定地点的行动，但无法自动证明全国的财产、人口、识字、信仰或动机。账本保留的证据说明是“起兵和主要诏令可靠；双方“奉天靖难”与“削藩”合法性措辞均具政治目的”。所以，官府标签、奏疏数字与后来的道德评语必须放回其文书目的，不能被写成不经分区核验的社会全貌。

这条事件更适合作为一组关系变化的锚点：中央方案需由地方官、书吏、士绅、宗族、商人与劳动者共同转译；不同家庭面对的税役、婚姻、迁徙和安全风险并不相等。材料只让我们看见其中有限的记录者。承认可见与不可见的界线，才能使制度史、文化史和区域史互相校正。

在账本条目\texttt{undefined}中，1400（建文二年）的“白沟河战役后燕军突破南下通道”发生于明中央军与燕军。本条把背景说明为建文军依赖李景隆等集中兵力围燕；燕军须夺取河北南下道路和补给，并以燕军获得向山东、直隶推进的机会；中央军统帅和兵源组织的弱点暴露概括可见影响。要理解它的社会后果，不能止于宫廷或战场：土地、赋役、司法、道路、性别化家庭劳动、城市市场、书院寺观和边海网络都会影响地方如何接收或改写制度。

本条依据《明太宗实录》建文二年条；《明史·成祖本纪》与《李景隆传》；北平、保定地方志战场条。这些来源能够较稳妥地支持所记年序、诏令、任命或特定地点的行动，但无法自动证明全国的财产、人口、识字、信仰或动机。账本保留的证据说明是“战役结局大体可靠；火器作用、军数和将领责任在双方叙事中夸张不一”。所以，官府标签、奏疏数字与后来的道德评语必须放回其文书目的，不能被写成不经分区核验的社会全貌。

这条事件更适合作为一组关系变化的锚点：中央方案需由地方官、书吏、士绅、宗族、商人与劳动者共同转译；不同家庭面对的税役、婚姻、迁徙和安全风险并不相等。材料只让我们看见其中有限的记录者。承认可见与不可见的界线，才能使制度史、文化史和区域史互相校正。

在账本条目\texttt{undefined}中，1401（建文三年）十二月的“东昌之战中燕军受挫，张玉战死”发生于明中央军与燕军。本条把背景说明为燕军南下后补给线拉长，建文军在山东仍可集结守军反击，并以燕军一度受阻并调整战略；持续内战增加沿线州县征发和交通损耗概括可见影响。要理解它的社会后果，不能止于宫廷或战场：土地、赋役、司法、道路、性别化家庭劳动、城市市场、书院寺观和边海网络都会影响地方如何接收或改写制度。

本条依据《明太宗实录》建文三年十二月条；《明史·张玉传》；山东地方志东昌战事条。这些来源能够较稳妥地支持所记年序、诏令、任命或特定地点的行动，但无法自动证明全国的财产、人口、识字、信仰或动机。账本保留的证据说明是“战役与张玉死亡可靠；胜败归因及伤亡数字不宜只用《太宗实录》”。所以，官府标签、奏疏数字与后来的道德评语必须放回其文书目的，不能被写成不经分区核验的社会全貌。

这条事件更适合作为一组关系变化的锚点：中央方案需由地方官、书吏、士绅、宗族、商人与劳动者共同转译；不同家庭面对的税役、婚姻、迁徙和安全风险并不相等。材料只让我们看见其中有限的记录者。承认可见与不可见的界线，才能使制度史、文化史和区域史互相校正。

在账本条目\texttt{undefined}中，1402（建文四年）六月的“燕军入南京，建文帝去向不明，朱棣即位”发生于明与燕军。本条把背景说明为燕军控制长江北岸与运河通道，建文朝内部指挥和防御体系瓦解，并以永乐朝开始重编前朝记录并重置人事；建文遗臣与继承合法性问题长期存在概括可见影响。要理解它的社会后果，不能止于宫廷或战场：土地、赋役、司法、道路、性别化家庭劳动、城市市场、书院寺观和边海网络都会影响地方如何接收或改写制度。

本条依据《明太宗实录》建文四年六月条；《明史·成祖本纪》《惠帝本纪》；南京明故宫考古与地方传说材料。这些来源能够较稳妥地支持所记年序、诏令、任命或特定地点的行动，但无法自动证明全国的财产、人口、识字、信仰或动机。账本保留的证据说明是“入城和即位可靠；宫城火灾、建文帝生死及出逃传说不能裁作确定事实”。所以，官府标签、奏疏数字与后来的道德评语必须放回其文书目的，不能被写成不经分区核验的社会全貌。

这条事件更适合作为一组关系变化的锚点：中央方案需由地方官、书吏、士绅、宗族、商人与劳动者共同转译；不同家庭面对的税役、婚姻、迁徙和安全风险并不相等。材料只让我们看见其中有限的记录者。承认可见与不可见的界线，才能使制度史、文化史和区域史互相校正。

在账本条目\texttt{undefined}中，1403（永乐元年）正月的“北平改称北京，永乐朝开始将其经营为北方政治中心”发生于明。本条把背景说明为永乐政权的军事基础在北平，且北元威胁要求皇帝更接近北方防线，并以北京获得新的行政与象征地位；大运河、宫城和卫戍工程成为迁都的前提概括可见影响。要理解它的社会后果，不能止于宫廷或战场：土地、赋役、司法、道路、性别化家庭劳动、城市市场、书院寺观和边海网络都会影响地方如何接收或改写制度。

本条依据《明太宗实录》永乐元年正月条；《明史·地理志》；北京城建考古调查资料。这些来源能够较稳妥地支持所记年序、诏令、任命或特定地点的行动，但无法自动证明全国的财产、人口、识字、信仰或动机。账本保留的证据说明是“改名和诏令可靠；改名不等于当年完成迁都或城建”。所以，官府标签、奏疏数字与后来的道德评语必须放回其文书目的，不能被写成不经分区核验的社会全貌。

这条事件更适合作为一组关系变化的锚点：中央方案需由地方官、书吏、士绅、宗族、商人与劳动者共同转译；不同家庭面对的税役、婚姻、迁徙和安全风险并不相等。材料只让我们看见其中有限的记录者。承认可见与不可见的界线，才能使制度史、文化史和区域史互相校正。

\subsection{环境、边防与海外连接}

在账本条目\texttt{undefined}中，1403（永乐元年）七月的“朝廷命解缙等编纂大型类书，后称《永乐大典》”发生于明。本条把背景说明为新政权需吸纳士人、整理经史与展示文化正统；内阁文臣获得组织工程的机会，并以形成跨门类文献汇编项目；其抄写、贮藏和后世散佚另构成一条物质史概括可见影响。要理解它的社会后果，不能止于宫廷或战场：土地、赋役、司法、道路、性别化家庭劳动、城市市场、书院寺观和边海网络都会影响地方如何接收或改写制度。

本条依据《明太宗实录》永乐元年条；《明史·解缙传》；《永乐大典》现存卷册题记。这些来源能够较稳妥地支持所记年序、诏令、任命或特定地点的行动，但无法自动证明全国的财产、人口、识字、信仰或动机。账本保留的证据说明是“诏修及主编群体可靠；成书规模和保存本不能直接说明知识的社会传播范围”。所以，官府标签、奏疏数字与后来的道德评语必须放回其文书目的，不能被写成不经分区核验的社会全貌。

这条事件更适合作为一组关系变化的锚点：中央方案需由地方官、书吏、士绅、宗族、商人与劳动者共同转译；不同家庭面对的税役、婚姻、迁徙和安全风险并不相等。材料只让我们看见其中有限的记录者。承认可见与不可见的界线，才能使制度史、文化史和区域史互相校正。

在账本条目\texttt{undefined}中，1404（永乐二年）的“朱高炽被立为皇太子，继承竞争暂获制度性裁定”发生于明。本条把背景说明为朱高炽、朱高煦各有不同的官僚和军功支持；永乐须固定继承秩序，并以太子在南京监国和文官网络中积累经验；兄弟竞争仍在永乐晚年反复出现概括可见影响。要理解它的社会后果，不能止于宫廷或战场：土地、赋役、司法、道路、性别化家庭劳动、城市市场、书院寺观和边海网络都会影响地方如何接收或改写制度。

本条依据《明太宗实录》永乐二年条；《明史·仁宗本纪》与《汉王高煦传》；杨士奇《东里文集》东宫奏议。这些来源能够较稳妥地支持所记年序、诏令、任命或特定地点的行动，但无法自动证明全国的财产、人口、识字、信仰或动机。账本保留的证据说明是“册立可靠；太子与汉王的个人动机不可由后出党争叙事简单化”。所以，官府标签、奏疏数字与后来的道德评语必须放回其文书目的，不能被写成不经分区核验的社会全貌。

这条事件更适合作为一组关系变化的锚点：中央方案需由地方官、书吏、士绅、宗族、商人与劳动者共同转译；不同家庭面对的税役、婚姻、迁徙和安全风险并不相等。材料只让我们看见其中有限的记录者。承认可见与不可见的界线，才能使制度史、文化史和区域史互相校正。

在账本条目\texttt{undefined}中，1405（永乐三年）六月的“郑和率舰队出航，开始第一次西洋航行”发生于明与印度洋港口政权。本条把背景说明为永乐朝借海上使团重建朝贡秩序并寻求海外情报和政治承认，并以官方舰队把东南港口、东南亚和印度洋使节网络连接起来；航行的财政和地方劳役成本需要另证概括可见影响。要理解它的社会后果，不能止于宫廷或战场：土地、赋役、司法、道路、性别化家庭劳动、城市市场、书院寺观和边海网络都会影响地方如何接收或改写制度。

本条依据《明太宗实录》永乐三年条；《明史·郑和传》；马欢《瀛涯胜览》相关航程。这些来源能够较稳妥地支持所记年序、诏令、任命或特定地点的行动，但无法自动证明全国的财产、人口、识字、信仰或动机。账本保留的证据说明是“出航可靠；船数、最大船型和贸易规模须区分官修记载、后出传说与物证”。所以，官府标签、奏疏数字与后来的道德评语必须放回其文书目的，不能被写成不经分区核验的社会全貌。

这条事件更适合作为一组关系变化的锚点：中央方案需由地方官、书吏、士绅、宗族、商人与劳动者共同转译；不同家庭面对的税役、婚姻、迁徙和安全风险并不相等。材料只让我们看见其中有限的记录者。承认可见与不可见的界线，才能使制度史、文化史和区域史互相校正。

在账本条目\texttt{undefined}中，1406（永乐四年）的“明以恢复陈氏为名出兵安南，进攻胡季犛政权”发生于明与大虞政权。本条把背景说明为胡氏取代陈朝引发流亡者求援；永乐朝将宗藩名义与南方扩张结合，并以战争迅速占领主要城市，却把明军置于持续的地方抵抗和补给压力之中概括可见影响。要理解它的社会后果，不能止于宫廷或战场：土地、赋役、司法、道路、性别化家庭劳动、城市市场、书院寺观和边海网络都会影响地方如何接收或改写制度。

本条依据《明太宗实录》永乐四年条；《明史·安南传》；越南《大越史记全书》相关纪年。这些来源能够较稳妥地支持所记年序、诏令、任命或特定地点的行动，但无法自动证明全国的财产、人口、识字、信仰或动机。账本保留的证据说明是“出兵和政治名义可靠；当地支持度与战场损失不能由明方檄文推断”。所以，官府标签、奏疏数字与后来的道德评语必须放回其文书目的，不能被写成不经分区核验的社会全貌。

这条事件更适合作为一组关系变化的锚点：中央方案需由地方官、书吏、士绅、宗族、商人与劳动者共同转译；不同家庭面对的税役、婚姻、迁徙和安全风险并不相等。材料只让我们看见其中有限的记录者。承认可见与不可见的界线，才能使制度史、文化史和区域史互相校正。

在账本条目\texttt{undefined}中，1407（永乐五年）的“明废大虞并设交趾承宣布政使司等机构”发生于明与交趾地区。本条把背景说明为明军俘获胡氏父子后选择直接统治而非扶立陈氏后裔，并以官署、卫所和税役被移植到交趾；高成本占领激化地方抵抗并牵制南方军力概括可见影响。要理解它的社会后果，不能止于宫廷或战场：土地、赋役、司法、道路、性别化家庭劳动、城市市场、书院寺观和边海网络都会影响地方如何接收或改写制度。

本条依据《明太宗实录》永乐五年条；《明史·地理志》与《安南传》；越南《大越史记全书》。这些来源能够较稳妥地支持所记年序、诏令、任命或特定地点的行动，但无法自动证明全国的财产、人口、识字、信仰或动机。账本保留的证据说明是“设治诏令可靠；行政建置不能证明红河三角洲和山区已被有效控制”。所以，官府标签、奏疏数字与后来的道德评语必须放回其文书目的，不能被写成不经分区核验的社会全貌。

这条事件更适合作为一组关系变化的锚点：中央方案需由地方官、书吏、士绅、宗族、商人与劳动者共同转译；不同家庭面对的税役、婚姻、迁徙和安全风险并不相等。材料只让我们看见其中有限的记录者。承认可见与不可见的界线，才能使制度史、文化史和区域史互相校正。

在账本条目\texttt{undefined}中，1407（永乐五年）的“郑和第二次航行出发，护送各国使节返航”发生于明与东南亚、南亚港口政权。本条把背景说明为第一次航行带回多国使者，明廷以返送维持官方互惠与海上声望，并以航行由一次远征转为可重复的使节循环；港口补给和翻译者作用更加关键概括可见影响。要理解它的社会后果，不能止于宫廷或战场：土地、赋役、司法、道路、性别化家庭劳动、城市市场、书院寺观和边海网络都会影响地方如何接收或改写制度。

本条依据《明太宗实录》永乐五年条；《明史·郑和传》；费信《星槎胜览》。这些来源能够较稳妥地支持所记年序、诏令、任命或特定地点的行动，但无法自动证明全国的财产、人口、识字、信仰或动机。账本保留的证据说明是“航行存在可靠；访问港口的先后和个别使团身份需与碑刻、海外材料互校”。所以，官府标签、奏疏数字与后来的道德评语必须放回其文书目的，不能被写成不经分区核验的社会全貌。

这条事件更适合作为一组关系变化的锚点：中央方案需由地方官、书吏、士绅、宗族、商人与劳动者共同转译；不同家庭面对的税役、婚姻、迁徙和安全风险并不相等。材料只让我们看见其中有限的记录者。承认可见与不可见的界线，才能使制度史、文化史和区域史互相校正。

在账本条目\texttt{undefined}中，1408（永乐六年）的“《永乐大典》进呈，类书编纂工程完成一阶段”发生于明。本条把背景说明为朝廷投入大批儒臣与抄手完成分类汇编，以文化工程巩固永乐正统，并以文献汇编成为内府收藏与后世辑佚的重要媒介；政治整肃中的编者命运也显示工程依赖宫廷权力概括可见影响。要理解它的社会后果，不能止于宫廷或战场：土地、赋役、司法、道路、性别化家庭劳动、城市市场、书院寺观和边海网络都会影响地方如何接收或改写制度。

本条依据《明太宗实录》永乐六年条；《明史·解缙传》；现存《永乐大典》卷端题记。这些来源能够较稳妥地支持所记年序、诏令、任命或特定地点的行动，但无法自动证明全国的财产、人口、识字、信仰或动机。账本保留的证据说明是“进呈可靠；“收尽天下书”是工程宣称，不能以之推定文献无遗漏”。所以，官府标签、奏疏数字与后来的道德评语必须放回其文书目的，不能被写成不经分区核验的社会全貌。

这条事件更适合作为一组关系变化的锚点：中央方案需由地方官、书吏、士绅、宗族、商人与劳动者共同转译；不同家庭面对的税役、婚姻、迁徙和安全风险并不相等。材料只让我们看见其中有限的记录者。承认可见与不可见的界线，才能使制度史、文化史和区域史互相校正。

在账本条目\texttt{undefined}中，1409（永乐七年）的“明设奴儿干都司，尝试经黑龙江下游联络和安置东北诸部”发生于明与奴儿干地区诸部。本条把背景说明为辽东经营需要北向水路、女真诸部关系和对北元残余的情报通道，并以明朝形成以朝贡、羁縻和使团为主的东北联系；永宁寺碑提供后续活动的物质锚点概括可见影响。要理解它的社会后果，不能止于宫廷或战场：土地、赋役、司法、道路、性别化家庭劳动、城市市场、书院寺观和边海网络都会影响地方如何接收或改写制度。

本条依据《明太宗实录》永乐七年条；《明史·地理志》与《奴儿干都司传》；永宁寺碑。这些来源能够较稳妥地支持所记年序、诏令、任命或特定地点的行动，但无法自动证明全国的财产、人口、识字、信仰或动机。账本保留的证据说明是“建置可靠；都司名义辖区不能等同于稳定领土控制”。所以，官府标签、奏疏数字与后来的道德评语必须放回其文书目的，不能被写成不经分区核验的社会全貌。

这条事件更适合作为一组关系变化的锚点：中央方案需由地方官、书吏、士绅、宗族、商人与劳动者共同转译；不同家庭面对的税役、婚姻、迁徙和安全风险并不相等。材料只让我们看见其中有限的记录者。承认可见与不可见的界线，才能使制度史、文化史和区域史互相校正。

在账本条目\texttt{undefined}中，1410（永乐八年）的“永乐帝亲征蒙古，北上击败本雅失里及相关部众”发生于明与鞑靼、瓦剌诸部。本条把背景说明为本雅失里处置明使及草原联盟变化挑战明廷北方威望，并以永乐以皇帝亲征重塑威慑；草原政治未被消除，明军远征补给负担持续扩大概括可见影响。要理解它的社会后果，不能止于宫廷或战场：土地、赋役、司法、道路、性别化家庭劳动、城市市场、书院寺观和边海网络都会影响地方如何接收或改写制度。

本条依据《明太宗实录》永乐八年条；《明史·成祖本纪》与《鞑靼传》；17世纪《蒙古源流》鞑靼、瓦剌叙事。这些来源能够较稳妥地支持所记年序、诏令、任命或特定地点的行动，但无法自动证明全国的财产、人口、识字、信仰或动机。账本保留的证据说明是“亲征和作战可靠；战果、敌军数字和草原部众归属须以多方材料审读”。所以，官府标签、奏疏数字与后来的道德评语必须放回其文书目的，不能被写成不经分区核验的社会全貌。

这条事件更适合作为一组关系变化的锚点：中央方案需由地方官、书吏、士绅、宗族、商人与劳动者共同转译；不同家庭面对的税役、婚姻、迁徙和安全风险并不相等。材料只让我们看见其中有限的记录者。承认可见与不可见的界线，才能使制度史、文化史和区域史互相校正。
